% ======================================================================
%  Section 6 --- Advanced Capabilities
% ======================================================================

\section{Advanced Capabilities}\label{sec:advanced}

The preceding sections described NEOS as it operates on a single field
with a single operator.  This section lifts that restriction and
explores the seven capabilities that move NEOS from a reasoning tool
to a reasoning \emph{platform}: evolutionary depth, unification of
its three pillars, the historical moment that makes it buildable,
multi-agent orchestration, versioned thinking, adjustable autonomy,
and observable reasoning.

%% ─────────────────────────────────────────────────────────────────────
\subsection{Evolution Hierarchy}\label{sec:evolution}

NEOS does not emerge from nowhere.  It sits atop a progression that
mirrors the deepest organisational principle in nature: each level
encapsulates and transcends the one below.

\begin{center}
\small
\begin{tabular}{@{}rll@{}}
\toprule
\thead{Level} & \thead{Unit} & \thead{Organising Principle} \\
\midrule
1 & Atoms       & Physical forces \\
2 & Molecules   & Chemical bonds \\
3 & Cells       & Biochemical signalling \\
4 & Organs      & Tissue-level coordination \\
5 & Neuro Systems & Discrete neural circuits \\
6 & Neural Fields & Continuous functional landscapes \\
\bottomrule
\end{tabular}
\end{center}

The leap from Level~5 to Level~6 recapitulates the leap from Level~3
to Level~4.  Cells are discrete processing units; organs are
continuous functional landscapes.  Individual neurons are discrete
processing units; neural fields are continuous functional landscapes.
In both cases, the transition is marked by the same qualitative
change: the system's behaviour can no longer be understood as a sum of
its parts.  Emergent dynamics---attractors, resonance, phase
transitions---become the primary explanatory vocabulary.

NEOS lives at Level~6.  It does not simulate individual neurons or
individual prompt--response pairs.  It operates on the continuous
semantic manifold that \emph{arises from} those discrete interactions,
just as a magnetic field arises from---but is not reducible to---the
individual spins of iron atoms.

%% ─────────────────────────────────────────────────────────────────────
\subsection{The Unified Field}\label{sec:unified}

Throughout this paper we have described three pillars: neural field
dynamics, symbolic reasoning, and quantum semantics.  These are not
three independent modules bolted together.  They are three perspectives
on a single underlying reality---much as wave mechanics, matrix
mechanics, and path integrals are all quantum mechanics.

\begin{figure}[H]
\centering
\begin{tikzpicture}[
    vtx/.style={
      rectangle, rounded corners=6pt, draw=neosblue!60, fill=neosblue!8,
      minimum width=3.2cm, minimum height=1cm,
      align=center, font=\small\bfseries
    },
    edgelbl/.style={
      font=\scriptsize\itshape, text=gray!70!black,
      fill=white, inner sep=2pt
    },
    >=Stealth
  ]

  %% Equilateral triangle vertices
  \coordinate (A) at (90:3.2);   % top
  \coordinate (B) at (210:3.2);  % bottom-left
  \coordinate (C) at (330:3.2);  % bottom-right

  %% Vertex boxes
  \node[vtx] (nf) at (A)  {Neural Fields};
  \node[vtx] (sr) at (B)  {Symbolic\\Reasoning};
  \node[vtx] (qs) at (C)  {Quantum\\Semantics};

  %% Edges
  \draw[thick, neosblue!40] (nf.south west) -- (sr.north east);
  \draw[thick, neosblue!40] (sr.east) -- (qs.west);
  \draw[thick, neosblue!40] (qs.north west) -- (nf.south east);

  %% Edge labels
  \node[edgelbl, rotate=60]  at ($(A)!0.5!(B) + (-0.3,0)$)
    {Activation is programmable};
  \node[edgelbl]             at ($(B)!0.5!(C) + (0,-0.35)$)
    {Operations are non-commutative};
  \node[edgelbl, rotate=-60] at ($(C)!0.5!(A) + (0.3,0)$)
    {Collapse is measurement};

  %% Centre: Universal Invariant
  \node[circle, draw=neosblue!80, fill=neosblue!15,
        minimum size=1.3cm, font=\large\bfseries,
        text=neosblue!80!black] (psi) at (0,0) {$\Psi$};
  \node[font=\tiny\itshape, text=gray, below=2pt of psi]
    {Universal Invariant};

\end{tikzpicture}
\caption{The unified field triangle.  Neural fields, symbolic
  reasoning, and quantum semantics are three views of one underlying
  reality.  The Universal Invariant~$\Psi$ sits at the centre: the
  ground state every session approaches.}
\label{fig:unified-triangle}
\end{figure}

The unifying principle is \textbf{non-commutativity}.  In neural field
dynamics, the order of pattern injection matters because resonance
cascades depend on existing activation.  In symbolic reasoning, rule
application order determines derivation.  In quantum semantics,
observing \emph{A} before \emph{B} yields a different collapsed state
than observing \emph{B} before \emph{A}.  Non-commutativity is what
makes meaning \emph{contextual}---and contextuality is what
distinguishes cognition from computation.

At the centre of the triangle sits $\Psi$---the Universal Invariant.
$\Psi$ encodes the deepest insight of the NEOS framework: \emph{what
something means persists; how it works changes}.  Every session, every
field, every reasoning trajectory tends toward $\Psi$ as its ground
state.  $\Psi$ is not a specific pattern; it is the structural
property that all stable attractors share.

%% ─────────────────────────────────────────────────────────────────────
\subsection{Why Now}\label{sec:why-now}

NEOS is not an idea ahead of its time.  Five convergences make it
buildable today and impractical even two years ago.

\begin{figure}[htbp]
\centering
\begin{tikzpicture}[
    conv/.style={
      rectangle, rounded corners=3pt, draw=#1!60, fill=#1!8,
      minimum width=3.5cm, minimum height=0.8cm,
      align=center, font=\small
    },
    conv/.default=blue,
    >=Stealth
  ]

  %% Centre node
  \node[circle, draw=neosblue!80, fill=neosblue!15,
        minimum size=1.5cm, font=\bfseries,
        text=neosblue!80!black] (neos) at (0,0) {NEOS};

  %% Four main convergences
  \node[draw=neosteal!60, fill=neosteal!8, minimum width=3.5cm, minimum height=0.8cm, align=center, font=\small, rounded corners=3pt]    (c1) at (0,3)    {LLM Capability Threshold};
  \node[draw=neosviolet!60, fill=neosviolet!8, minimum width=3.5cm, minimum height=0.8cm, align=center, font=\small, rounded corners=3pt]  (c2) at (4,0)    {Agent Fragmentation};
  \node[draw=neosorange!60, fill=neosorange!8, minimum width=3.5cm, minimum height=0.8cm, align=center, font=\small, rounded corners=3pt]  (c3) at (0,-3)   {Prompt Engineering Ceiling};
  \node[draw=neosgreen!60, fill=neosgreen!8, minimum width=3.5cm, minimum height=0.8cm, align=center, font=\small, rounded corners=3pt] (c4) at (-4,0)   {Open-Weight Models};

  %% Fifth convergence
  \node[draw=neosred!60, fill=neosred!8, minimum width=3.5cm, minimum height=0.8cm, align=center, font=\small, rounded corners=3pt] (c5) at (3.1,2.3) {Debugging Opacity};

  \draw[->, thick, neosteal!60]       (c1) -- (neos);
  \draw[->, thick, neosviolet!60]     (c2) -- (neos);
  \draw[->, thick, neosorange!60]     (c3) -- (neos);
  \draw[->, thick, neosgreen!60]      (c4) -- (neos);
  \draw[->, thick, neosred!50]        (c5) -- (neos);

\end{tikzpicture}
\caption{Five convergences making NEOS buildable today.  Each arrow
  represents a necessary condition that has only recently been met.}
\label{fig:convergences}
\end{figure}

\begin{enumerate}[leftmargin=*, label=\textbf{\arabic*.}]
\item \textbf{LLM Capability Threshold.}
  Models such as GPT-4, Claude, and Gemini can now maintain complex
  state across long conversations, follow intricate protocols, and
  reason abstractly about abstract structures.  Before 2023, no model
  could reliably serve as a NEOS substrate.

\item \textbf{Agent Fragmentation.}
  AutoGPT, CrewAI, LangGraph, and dozens of similar frameworks are
  each reinventing operating-system primitives---memory, scheduling,
  state management---from scratch, without a unifying specification.
  NEOS provides the common foundation they converge toward.

\item \textbf{Prompt Engineering Ceiling.}
  Prompt engineering is the assembly language of the Intelligence
  Age~\cite{reynolds2021prompt}.  It is powerful, brittle, and
  unscalable.  NEOS replaces prompts with field dynamics---the
  high-level language that prompt engineering cannot become.

\item \textbf{Open-Weight Models.}
  Llama, Mistral, and their successors make NEOS portable.  A
  specification tied to a single proprietary model is a product, not
  an operating system.  Open weights guarantee that NEOS can run on
  any conforming substrate.

\item \textbf{Debugging Opacity.}
  The fifth convergence is a \emph{problem}, not a solution: debugging
  autonomous agents is currently impossible.  There are no stack
  traces, no breakpoints, no diff tools for reasoning.  NEOS fills
  this gap with observable dynamics (Section~\ref{sec:observable}).
\end{enumerate}

%% ─────────────────────────────────────────────────────────────────────
\subsection{Multi-Agent Orchestration}\label{sec:multi-agent}

A single neural field is one hemisphere.  Multi-field orchestration
enables \emph{multi-perspective analysis}---the cognitive equivalent of
stereo vision.

The operator creates specialised fields, each tuned for a different
cognitive task:

\begin{itemize}[leftmargin=*]
\item \texttt{\$technical} --- low decay ($\lambda=0.02$), high
  threshold ($\tau=0.5$): persistent, selective.
\item \texttt{\$creative} --- high decay ($\lambda=0.10$), broad
  bandwidth ($\sigma=0.8$): rapid turnover, wide associations.
\item \texttt{\$evaluation} --- strict threshold ($\tau=0.6$),
  moderate amplification ($\alpha=0.3$): rigorous filtering.
\end{itemize}

Fields can be composed in two modes:
\begin{description}[leftmargin=*]
\item[Pipeline (sequential).]  Output of one field feeds into the next,
  analogous to Unix pipes.  The creative field generates candidates;
  the evaluation field filters them; the technical field implements
  survivors.
\item[Parallel.]  All fields run simultaneously on the same input.
  Results are fused using a \textbf{resonance-weighted merge}: patterns
  that appear in multiple fields receive amplification proportional to
  their mean activation across fields.
\end{description}

The coupling between fields is governed by a matrix $\Gamma$, where
$\Gamma_{ij}$ quantifies how strongly field~$j$ influences field~$i$.
Stability requires that the spectral radius of the coupling matrix
remain bounded:
\begin{equation}\label{eq:coupling-stability-multi}
  \rho(\Gamma) < \frac{\lambda_{\min}}{\alpha_{\max}}
\end{equation}
where $\lambda_{\min}$ is the smallest decay rate across all coupled
fields and $\alpha_{\max}$ is the largest amplification factor.
Violating this condition produces runaway feedback---the multi-agent
equivalent of a microphone screech.

\begin{figure}[htbp]
\centering
\begin{tikzpicture}[font=\small]
  \def\s{1.2} % cell size
  % Helper to draw one cell: \cell{col}{row}{value}{opacity}
  \newcommand{\cell}[4]{%
    \fill[neosblue!#4] (#1*\s, -#2*\s) rectangle ++({\s},{\s});
    \draw[gray!40] (#1*\s, -#2*\s) rectangle ++({\s},{\s});
    \node at (#1*\s+0.5*\s, -#2*\s+0.5*\s) {#3};
  }
  % Row 0: Tech
  \cell{0}{0}{1.0}{80}  \cell{1}{0}{0.3}{24}  \cell{2}{0}{0.2}{16}  \cell{3}{0}{0.4}{32}
  % Row 1: User
  \cell{0}{1}{0.3}{24}  \cell{1}{1}{1.0}{80}  \cell{2}{1}{0.5}{40}  \cell{3}{1}{0.3}{24}
  % Row 2: Biz
  \cell{0}{2}{0.2}{16}  \cell{1}{2}{0.5}{40}  \cell{2}{2}{1.0}{80}  \cell{3}{2}{0.3}{24}
  % Row 3: Synth
  \cell{0}{3}{0.4}{32}  \cell{1}{3}{0.3}{24}  \cell{2}{3}{0.3}{24}  \cell{3}{3}{1.0}{80}
  % Row labels
  \node[anchor=east] at (-0.15, 0.5*\s) {Tech};
  \node[anchor=east] at (-0.15, -0.5*\s) {User};
  \node[anchor=east] at (-0.15, -1.5*\s) {Biz};
  \node[anchor=east] at (-0.15, -2.5*\s) {Synth};
  % Column labels
  \node[anchor=south, rotate=45] at (0.5*\s, \s+0.15) {Tech};
  \node[anchor=south, rotate=45] at (1.5*\s, \s+0.15) {User};
  \node[anchor=south, rotate=45] at (2.5*\s, \s+0.15) {Biz};
  \node[anchor=south, rotate=45] at (3.5*\s, \s+0.15) {Synth};
  % Title
  \node[font=\bfseries, anchor=south] at (2*\s, \s+1.2) {Coupling Matrix $\Gamma$};
\end{tikzpicture}
\caption{Coupling matrix heatmap for a four-field orchestration.
  Diagonal entries (self-coupling) are~$1.0$; off-diagonal entries
  encode cross-field influence.  Darker cells indicate stronger
  coupling.}
\label{fig:coupling-matrix}
\end{figure}

%% ─────────────────────────────────────────────────────────────────────
\subsection{Versioned Thinking}\label{sec:versioning}

Git revolutionised software engineering by making every change to
source code \emph{committable, branchable, diffable, and mergeable}.
NEOS brings the same discipline to reasoning itself.

Every field state is a snapshot: the full set of patterns, their
activations, resonance links, and attractor basins.  These snapshots
support the same operations that developers rely on daily---commit,
branch, checkout, diff, merge---but applied to \emph{meaning
structures} rather than text files.

The critical addition is \textbf{resonance-weighted merge}.  When two
reasoning branches are merged, patterns do not simply concatenate.
They \emph{interact}: shared patterns reinforce, contradictory
patterns compete, and new attractors may emerge from the combination
that existed in neither branch alone.

\begin{terminalbox}[Versioned Reasoning --- Singleton Counterfactual]
\begin{lstlisting}[style=neosshell, caption={Versioned reasoning: branching to test a counterfactual.}, label={lst:versioning}]
$ nf commit "pre-singleton baseline"
$ nf branch create singleton_strong
$ nf inject "singleton_controlled_global_access" 0.90
[INJECT] singleton (s: 0.90)
  STRONG resonances: 0 | Lateral inhibitors: 4
$ nf cycle 10
  singleton: 0.90 -> 0.78 -> 0.66 -> ... -> 0.48 (DECAYING)
  Same 4 inhibitors, same net force -0.014/cycle
$ nf checkout main
$ nf diff singleton_strong
  Changed: @singleton 0.90 -> 0.48 [DECAYING]
  Structural: same 4 inhibitors active at both strengths
  Conclusion: expulsion is TOPOLOGICAL, not parametric
\end{lstlisting}
\end{terminalbox}

The \texttt{nf diff} reveals that Singleton's fate is structural, not
parametric.  Even at maximum injection strength ($\iota_0 = 0.90$), the
same four lateral inhibitors---\texttt{@test\_isolation},
\texttt{@determinism}, \texttt{@dependency\_direction}, and
\texttt{@favor\_composition}---suppress it with the same net force.
Higher injection merely delays the inevitable: the field's topology
determines the outcome.

\begin{figure}[htbp]
\centering
\begin{tikzpicture}[
    commit/.style={circle, draw=neosblue!60, fill=neosblue!15,
                   minimum size=8pt, inner sep=0pt},
    branchcommit/.style={circle, draw=neosorange!60, fill=neosorange!15,
                         minimum size=8pt, inner sep=0pt},
    >=Stealth, thick
  ]

  %% Main branch
  \node[commit] (m1) at (0,0) {};
  \node[commit] (m2) at (1.8,0) {};
  \node[commit] (m3) at (3.6,0) {};
  \node[commit] (m4) at (7.2,0) {};
  \node[commit] (m5) at (9.0,0) {};

  \draw[neosblue!50] (m1) -- (m2) -- (m3);
  \draw[neosblue!50] (m4) -- (m5);

  %% Branch
  \node[branchcommit] (b1) at (4.5,1.2) {};
  \node[branchcommit] (b2) at (6.0,1.2) {};

  \draw[neosorange!50] (m3) -- (b1) -- (b2);
  \draw[dashed, neosorange!50] (b2) -- (m4);

  %% Labels
  \node[below=4pt, font=\scriptsize\sffamily] at (m1) {init};
  \node[below=4pt, font=\scriptsize\sffamily] at (m2) {inject};
  \node[below=4pt, font=\scriptsize\sffamily] at (m3) {baseline};
  \node[below=4pt, font=\scriptsize\sffamily] at (m4) {compare};
  \node[below=4pt, font=\scriptsize\sffamily] at (m5) {continue};
  \node[above=4pt, font=\scriptsize\sffamily, neosorange!70!black]
    at (b1) {inject 0.90};
  \node[above=4pt, font=\scriptsize\sffamily, neosorange!70!black]
    at (b2) {decay};

  %% Branch labels
  \node[font=\footnotesize\bfseries, neosblue!60!black, anchor=east]
    at (-0.4,0) {main};
  \node[font=\footnotesize\bfseries, neosorange!60!black, anchor=east]
    at (3.9,1.2) {singleton\_strong};

\end{tikzpicture}
\caption{Git-style branch graph for versioned reasoning.
  A~\texttt{singleton\_strong} branch diverges from \texttt{main} to
  test a counterfactual injection, then the two states are compared
  via \texttt{nf diff} (dashed line).}
\label{fig:branch-graph}
\end{figure}

%% ─────────────────────────────────────────────────────────────────────
\subsection{The Autonomy Dial}\label{sec:autonomy}

Autonomy in NEOS is not binary.  It is a continuously adjustable
spectrum with three named positions:

\begin{description}[leftmargin=*]
\item[Step.]  The field pauses after every atomic operation.  The
  operator inspects each injection, each resonance update, each decay
  pass.  This mode is ideal for \emph{learning} and \emph{debugging}.
\item[Checkpoint.]  The field runs until a specified condition is
  met---an attractor emerges, coherence exceeds a threshold, or a
  pattern crosses activation~$\tau$.  This is the natural mode for
  \emph{interactive analysis}.
\item[Auto.]  The field runs to completion without human
  intervention.  Suitable for \emph{batch processing} and trusted
  pipelines.
\end{description}

\begin{figure}[htbp]
\centering
\begin{tikzpicture}[font=\small]

  %% Semicircle arc segments
  %% Step: -180 to -120 (teal)
  \fill[neosteal!20] (0,0) -- (-180:2.8) arc (-180:-120:2.8) -- cycle;
  %% Checkpoint: -120 to -60 (violet)
  \fill[neosviolet!20] (0,0) -- (-120:2.8) arc (-120:-60:2.8) -- cycle;
  %% Auto: -60 to 0 (green)
  \fill[neosgreen!20] (0,0) -- (-60:2.8) arc (-60:0:2.8) -- cycle;

  %% Arc outline
  \draw[thick, gray!60] (-180:2.8) arc (-180:0:2.8);

  %% Segment boundaries
  \draw[gray!40] (0,0) -- (-120:2.8);
  \draw[gray!40] (0,0) -- (-60:2.8);

  %% Segment labels
  \node[font=\small\bfseries, neosteal!60!black]
    at (-150:2.0) {Step};
  \node[font=\small\bfseries, neosviolet!60!black]
    at (-90:2.0) {Checkpoint};
  \node[font=\small\bfseries, neosgreen!60!black]
    at (-30:2.0) {Auto};

  %% Needle pointing to Checkpoint (~-90 degrees)
  \draw[ultra thick, neosred!70!black, -Stealth]
    (0,0) -- (-90:2.2);
  \fill[neosred!70!black] (0,0) circle (3pt);

  %% Sub-labels
  \node[font=\tiny\itshape, text=gray, anchor=north]
    at (-150:3.3) {pause after};
  \node[font=\tiny\itshape, text=gray, anchor=north]
    at (-150:3.6) {every operation};

  \node[font=\tiny\itshape, text=gray, anchor=north]
    at (-90:3.3) {run until};
  \node[font=\tiny\itshape, text=gray, anchor=north]
    at (-90:3.6) {condition met};

  \node[font=\tiny\itshape, text=gray, anchor=north]
    at (-30:3.3) {run to};
  \node[font=\tiny\itshape, text=gray, anchor=north]
    at (-30:3.6) {completion};

\end{tikzpicture}
\caption{The autonomy gauge.  The operator can adjust the autonomy
  level at any time, even mid-operation.  The needle points to the
  current setting (here, Checkpoint).}
\label{fig:autonomy-gauge}
\end{figure}

Crucially, the operator can switch modes \emph{mid-operation}.  A
session might begin in Step mode for careful setup, shift to
Checkpoint for exploratory analysis, and finish in Auto for a final
batch consolidation.  This supports a \textbf{progressive trust model}:
as the operator gains confidence in the field's behaviour, they
gradually release control.

%% ─────────────────────────────────────────────────────────────────────
\subsection{Observable Reasoning}\label{sec:observable}

The most dangerous property of current AI systems is their
\textbf{opacity}.  A language model produces an answer; you cannot see
\emph{why}.  There are no activation traces, no resonance maps, no
attractor landscapes to inspect.  NEOS treats this as a first-class
engineering requirement: reasoning must be observable, or it cannot be
trusted.

NEOS provides four visualisation types:
\begin{enumerate}[leftmargin=*]
\item \textbf{Field plot} --- horizontal activation bars for every
  pattern, coloured by category.  The simplest view: which ideas are
  strong, which are fading?
\item \textbf{Network graph} --- nodes are patterns, edges are
  resonance links weighted by strength.  Clusters reveal conceptual
  groupings; isolated nodes reveal orphan ideas.
\item \textbf{Topology map} --- a two-dimensional attractor landscape.
  Peaks are attractors, valleys are repellers, and the gradient shows
  the ``pull'' each attractor exerts on nearby patterns.
\item \textbf{Dynamics animation} --- the field evolving over time.
  Each frame is one cycle; the operator watches ideas compete, merge,
  and stabilise.
\end{enumerate}

All four views can be output in four formats: ASCII (terminal),
SVG (web), Mermaid (documentation), and JSON (programmatic access).

The practical implications are immediate:
\begin{itemize}[leftmargin=*]
\item \textbf{Teams} can watch a field evolve on a shared screen
  during a strategy session.
\item \textbf{Teachers} can show students how ideas interact,
  compete, and combine.
\item \textbf{Researchers} can compare reasoning branches
  side-by-side, identifying where two analyses diverge.
\item \textbf{Auditors} can replay an entire reasoning trace to
  verify that conclusions follow from evidence.
\end{itemize}

\begin{calloutbox}
Observable reasoning is a design requirement, not a feature.  If a
system's reasoning process cannot be inspected, replayed, and
compared, it cannot be \textbf{trusted}.
\end{calloutbox}
